
Nuclear power has supplied the United States with large quantities of
low-carbon electricity for more than six decades. The fuel discharged from
those reactors, however, remains a long term national responsibility. By the
end of 2025, commercial reactors in the United States had permanently
discharged roughly \(95{,}000\) \gls{MTHM} of \gls{UNF}, and the
inventory continues to grow by approximately \(2{,}000\) \gls{MTHM} each year
\cite{inside_2025,carter_banerjee_inventory_2024}. Extrapolating this rate places the
inventory at approximately \(99{,}000\) \gls{MTHM} in 2026. This material is distributed across
more than one hundred sites in thirty-nine states, including operating nuclear
power plants, decommissioned reactor sites, and independent spent fuel storage
installations. Because the United States does not yet operate a permanent
geologic repository for commercial spent fuel, most of the inventory remains
stored at or near the sites where it was generated
\cite{nrc_spent_fuel_storage}.

The challenge is not simply that the inventory is large. It is also highly
heterogeneous. The phrase ``the U.S. used nuclear fuel inventory'' can suggest
a single stockpile with one representative composition, but the actual
inventory is the accumulated result of decades of reactor operation.
Assemblies discharged in different years were fabricated with different
initial enrichments, irradiated to different burnups, operated in different
reactor types, and cooled for different lengths of time. Their remaining
heavy metal masses and isotopic compositions therefore vary considerably.
Some assemblies contain larger quantities of residual fissile material,
plutonium, or minor actinides than others. Some have cooled for several
decades, while others were discharged recently and continue to generate
substantial decay heat. Any strategy that attempts to process, recycle, or
transmute this material must eventually confront those differences.

\GLS{UNF} also changes with time. During the first decades after
discharge, its heat generation and radiation field are strongly influenced by
fission products. As those nuclides decay, the relative importance of
plutonium and the minor actinides increases. Their long half-lives contribute
to the need to isolate untreated fuel over timescales on the order of
\(10^{5}\) years \cite{arpa_e_newton_2025}. Among the long-lived components
of used fuel, the actinides are distinctive because they can be destroyed
through fission rather than managed only through continued storage and
radioactive decay. That possibility does not eliminate the need for a
repository, nor does it remove every radioactive waste stream. It does,
however, create the possibility of changing the quantity and composition of
the material that must ultimately be isolated.

\Glspl{ADS} have been proposed as one technology for carrying
out this transmutation. An \gls{ADS} couples an externally generated neutron
source to a subcritical multiplying core. Because the core remains below
criticality, its operation depends on the external source, allowing fuel
compositions with high concentrations of minor actinides to be considered
under conditions that would be more difficult to accommodate in a critical
reactor. Detailed studies have examined accelerator reliability, spallation
targets, subcritical core physics, liquid metal coolants, fuel behavior, and
the rates at which plutonium and minor actinides could be consumed. These
studies establish the physical basis for accelerator driven transmutation and
provide increasingly realistic designs for future facilities.

A technically feasible core, however, is not the same as a viable
transmutation strategy. A reactor physics calculation may show how much
actinide material one system can consume, but it does not determine where that
material will come from, which assemblies will be processed first, how the
feed will change over time, or how many facilities will be required. Before an
\gls{ADS} can consume actinides, spent fuel must be selected from storage,
transported, cooled as needed, separated, and converted into a suitable feed
form. The transmutation facility must then be deployed, supplied, operated,
and eventually integrated with storage, recycle, and waste management
facilities. The material discharged from the \gls{ADS} may require further
processing or repeated recycle before the desired degree of actinide
destruction is achieved.

The central problem is therefore a fuel cycle problem rather than only a
single core problem. It requires a model that can follow material from the
existing inventory, through selection and processing, into a transmutation
facility, and finally into the inventories and waste streams that remain.
Such a model must operate over timescales of years or decades while retaining
enough information to distinguish among assemblies. It must also allow
facilities to compete for resources, enforce throughput limits, and respond
to changes in material availability. These capabilities are difficult to
represent using a single average spent fuel composition.

Many previous fuel cycle studies have necessarily simplified the inventory by
representing it as an aggregated material stream or a small set of
representative recipes \cite{rodriguez_nuclear_2020,nea_regional_scenarios_2009,coquelet_pascal_cosi6_2015}.
That approach is useful when comparing generic
future fuel cycles or studying equilibrium reactor fleets. It is less
satisfactory when the starting point is the existing U.S. inventory. A
national average recipe cannot identify whether the material it represents
comes from older or newer assemblies, whether it has high or low burnup,
whether it is concentrated at a few sites or dispersed across the country, or
which assemblies would actually be selected under a particular management
strategy. These differences may affect the amount and quality of material
available for transmutation and the characteristics of the inventory that
remains untreated.

Recent inventory databases make a more realistic representation possible.
The \gls{STANDARDS}~5.0.1 database \cite{peterson_unf-st&dards_2017,peterson_additional_2017}
contains assembly-level information describing
the U.S. commercial spent fuel inventory, including attributes such as
discharge date, burnup, initial enrichment, mass, and calculated isotopic
composition. This level of detail creates an opportunity to treat the
national inventory not as a single homogeneous stockpile, but as a collection
of discrete resources that can be selected according to explicit criteria.
The database alone, however, does not describe how those assemblies would
move through a future fuel cycle. The information must be incorporated into a
dynamic simulator capable of representing transactions among facilities.

The work is conducted within the Advanced Reactors and Fuel Cycles group at
\gls{UIUC}, where agent-based fuel cycle simulation is used to examine the
system level consequences of advanced nuclear technologies. It also
contributes to the UIUC led \gls{EINSTEIN} project under the \gls{ARPAE} \gls{NEWTON} program,
which brings together neutronics, nuclear data, materials, fuel cycle, and
techno-economic analyses to evaluate externally driven transmutation systems.
The topic of this thesis involves simulation with \Cyclus,
detailed spent fuel inventory data, and higher fidelity reactor calculations,
to fill a gap between the available inventory information and
existing system level transmutation models.
This thesis connects the fuel assemblies that exist today with the
transmutation technologies being proposed for the future.

It uses the \Cyclus nuclear fuel cycle simulation framework to make
that connection. \Cyclus represents the fuel cycle as a system of interacting
agents, including regions, institutions, and facilities \cite{huff_fundamental_2016}. Individual
facilities request, offer, and exchange material through the Dynamic Resource
Exchange \cite{giddenMethodologyDeterminingDynamic2016},
while modular software components known as archetypes define their
behavior. This structure makes it possible to represent the U.S. inventory as
an active supplier rather than as a static input file \cite{bae_deep_2020}. It also allows that
supplier to interact with processing, storage, and transmutation facilities
within the same simulation.

Two new facility archetypes are developed for this purpose. The
\texttt{us\_inventory} archetype represents individual U.S. spent fuel
assemblies as selectable material resources. Each assembly retains metadata
including its heavy metal mass, discharge date, burnup, initial enrichment,
and isotopic composition derived from \gls{STANDARDS}~5.0.1.
Rather than collapsing these records into one average composition, the facility supplies
material through the \Cyclus Dynamic Resource Exchange according to
user defined throughput limits and inventory selection policies.

The second archetype \texttt{ads}, represents an \gls{ADS}. Its material transformation is
based on precomputed Serpent depletion calculations, allowing the
fuel cycle simulation to incorporate higher fidelity neutronics results
without performing a full transport and depletion calculation during every
simulation time step. Coupling the two archetypes makes it possible to follow
material from individual inventory records into the modeled transmutation
system and to observe how the remaining inventory changes as material is
withdrawn.

The selection policies considered in this work rank available assemblies
according to characteristics such as discharge date, burnup, and initial
enrichment. These policies are not intended to represent final operational
decisions or recommendations for U.S. waste management. Instead, they provide
controlled cases for examining how different assumptions about access to the
inventory affect the material supplied to a transmutation system. For
example, prioritizing older assemblies may produce a different feed
composition than prioritizing high burnup or high enrichment assemblies, even
when the same total mass is supplied. The simulation records these
differences and preserves the connection between the supplied material and
the assemblies from which it originated.

The primary contribution of this thesis is a material flow model rather than
a new accelerator driven reactor design or a full evaluation of whether the
United States should deploy accelerator driven transmutation. A decision of
that kind would still require detailed neutronics, thermal hydraulics,
materials, safety, economic, transportation, safeguards, and repository
analyses. What this work provides instead is a model that connects a
realistic, assembly resolved representation of the U.S. inventory to an
\gls{ADS} facility inside a dynamic fuel cycle simulation.

This connection provides a foundation for broader policy and deployment
studies. A future analysis could add site specific transportation,
separations capacity, construction schedules, facility siting, operating
costs, or alternative transmutation technologies. Those questions cannot be
fully grounded in the present inventory, however, unless the model first
represents the assemblies that actually exist and the rules by which they
become available. By preserving assembly-level differences, this work makes
the assumptions behind a transmutation scenario more transparent and its
results more traceable.

The broader motivation is to ground the discussion of transmutation in how a
real national inventory could supply such facilities over time, rather than in
what a single idealized facility can destroy on its own. This emphasis matters
because the performance of a technology in isolation may differ from its
performance as part of an operating fuel cycle. Through the integration of
assembly-level inventory data and an \gls{ADS} representation in \Cyclus, this
thesis provides one step toward evaluating that larger system.

The remainder of this thesis is organized as follows.
Chapter~\ref{ch:background} reviews \gls{UNF} management in the
United States, the \Cyclus framework, accelerator driven systems, the
\gls{EINSTEIN} program, and prior transmutation fuel cycle simulations. It concludes
by identifying the research gap addressed in this work.
Chapter~\ref{ch:methodology} describes the design and implementation of the
\texttt{us\_inventory} and \texttt{ads} archetypes, the preparation of their
input data, and the simulation scenarios considered.
Chapter~\ref{ch:results} presents the simulation results and compares the
effects of the alternative inventory selection policies.
Finally, Chapter~\ref{ch:conclusion} discusses the implications and
limitations of the results and identifies directions for future development.